% Generated draft: requires mathematical and source review.
\documentclass[11pt]{article}
\usepackage[margin=25mm]{geometry}
\usepackage{fontspec}
\setmainfont{Latin Modern Roman}
\usepackage{amsmath,amssymb,mathtools,unicode-math}
\setmathfont{Latin Modern Math}
\usepackage{enumitem,needspace,xurl,hyperref}
\setlength{\parindent}{0pt}
\setlength{\parskip}{5pt}
\setlength{\emergencystretch}{3em}
\newcommand{\N}{\mathbb{N}}
\newcommand{\Z}{\mathbb{Z}}
\newcommand{\Q}{\mathbb{Q}}
\newcommand{\R}{\mathbb{R}}
\newcommand{\C}{\mathbb{C}}
\newcommand{\F}{\mathbb{F}}
\newcommand{\A}{\mathbb{A}}
\newcommand{\PP}{\mathbb{P}}
\newcommand{\E}{\mathbb{E}}
\newcommand{\eps}{\varepsilon}
\newcommand{\dd}{\,\mathrm{d}}
\begin{document}
\section*{UnsolvedMath ID: 30001331}
\textbf{Source problem code:} catalogue entry
\subsection{Definitions and mathematical statement}
For an abstract group \(\Gamma\) and a compact Hausdorff space \(X\), interpret \(C(X)\) provisionally as the unital ring \(C(X,\C)\) of continuous complex-valued functions, with pointwise operations. The supplied record does not explicitly specify the scalar field; the formulation below uses this usual convention and requires confirmation against the source. The algebraic group ring \(C(X)[\Gamma]\) consists of finite sums \(\sum_{\gamma\in\Gamma}a_\gamma\gamma\), with coefficients in \(C(X)\), coefficientwise addition, and multiplication determined by the group law and by coefficients commuting with group elements. No completion is taken.

For a unital ring \(A\), let \(K_0^{\mathrm{alg}}(A)\) be the Grothendieck group of isomorphism classes of finitely generated projective right \(A\)-modules, with relations \([P\oplus Q]=[P]+[Q]\). A unital ring homomorphism induces a homomorphism on \(K_0^{\mathrm{alg}}\) by extension of scalars. For the empty space, use \(C(\varnothing)=0\) and \(K_0^{\mathrm{alg}}(0)=0\).

A continuous map \(f\colon X\to Y\) induces the pullback ring homomorphism
\[
 f^*\colon C(Y)[\Gamma]\longrightarrow C(X)[\Gamma],\qquad
 \sum_\gamma a_\gamma\gamma\longmapsto
 \sum_\gamma(a_\gamma\circ f)\gamma.
\]
Write \(F_\Gamma(X)=K_0^{\mathrm{alg}}(C(X)[\Gamma])\) and \(F_\Gamma(f)=K_0^{\mathrm{alg}}(f^*)\); thus \(F_\Gamma\) is contravariant.

\textbf{Conjecture, conditional on the scalar convention above.} For every group \(\Gamma\), every pair of compact Hausdorff spaces \(X,Y\), and every pair of continuous maps \(f_0,f_1\colon X\to Y\) admitting a continuous homotopy \(H\colon X\times[0,1]\to Y\) with \(H(x,0)=f_0(x)\) and \(H(x,1)=f_1(x)\) for all \(x\in X\), one has
\[
 F_\Gamma(f_0)=F_\Gamma(f_1)\colon F_\Gamma(Y)\longrightarrow F_\Gamma(X).
\]
Equivalently, for every group \(\Gamma\) and compact Hausdorff space \(X\), the projection \(p\colon X\times[0,1]\to X\) induces an isomorphism
\[
 K_0^{\mathrm{alg}}(p^*)\colon
 K_0^{\mathrm{alg}}(C(X)[\Gamma])
 \xrightarrow{\ \cong\ }
 K_0^{\mathrm{alg}}(C(X\times[0,1])[\Gamma]).
\]
There are no torsion, finite-generation, or polyhedral restrictions in the supplied conjecture. The supplied partial-results assessment does not specify enough hypotheses to formulate its special cases precisely.
\subsection{Short English statement}
For every group \(\Gamma\), ask whether homotopic maps between compact Hausdorff spaces induce the same map on the algebraic \(K_0\)-groups of their function group rings. Equivalently, adjoining an interval factor should induce an isomorphism. Here \(C(X)\) is provisionally interpreted as complex-valued continuous functions; the scalar convention needs source confirmation.
\subsection{Sources}
\begin{itemize}
\item Supplied local UnsolvedMath registry record, permanent ID \(30001331\), source problem code \texttt{catalogue entry}, entitled \emph{Homotopy Invariance of Algebraic K-Theory for Group Rings}, in the Oberwolfach Reports Open Problems collection. No individual proposer is supplied; the proposed year is recorded as 2009.
\item Supplied report citation: \emph{Noncommutative Geometry (2010)}, Oberwolfach Reports, pp.\ 2271--2334, DOI \texttt{10.4171/owr/2009/41}. Source URL: \url{https://doi.org/10.4171/owr/2009/41}. The background describes Oberwolfach Reports as records of mathematical workshops and their open-problem discussions.
\item Supplied literature source, labeled ``original Oberwolfach report 4084'': \url{https://doi.org/10.4171/owr/2009/41}.
\item Literature assessment dated \texttt{2026-08-21}: supplied dated curation, not independently verified here. It labels the problem partially solved and reports proofs for important classes, including torsion-free groups under a Farrell--Jones hypothesis and a polyhedral setting, while reporting that the assertion for every group remains open. The supplied background explicitly requests independent verification of time-sensitive status and cited claims.
\item The registry supplies the statement-status label \texttt{corrected\_verified} and says that the mathematical content was checked against the official report, with notation, context, and question form normalized without changing the source-level claim. This is a supplied verification assertion, not an independent verification performed for this draft.
\end{itemize}
\end{document}
